\chapter{Methodology}

\section{System Design}

\subsection{Embedded Subsystem (ESP32)}

\subsubsection{Heart Rate}

Heart rate is derived from the photoplethysmography (PPG) infrared signal by detecting
peaks in the waveform — each peak corresponding to one cardiac cycle.

\subsubsection{SpO\textsubscript{2}}

Blood oxygen saturation is estimated by pulse oximetry, exploiting the difference in
optical absorption between oxyhaemoglobin (HbO\textsubscript{2}) and deoxyhaemoglobin
(Hb) at two wavelengths.

\subsubsection{Delta-Threshold Upload Filter}

Records are enqueued only when the measured value has changed by more than a defined
threshold since the last successful push.
Let $x_n$ be the current reading and $\hat{x}$ the last transmitted value.
The upload condition for room data is:

\begin{equation}
	\text{Upload}_{\text{room}} \;\Leftrightarrow\;
	\left| T_n - \hat{T} \right| > \delta_T \;\lor\;
	\left| H_n - \hat{H} \right| > \delta_H
	\label{eq:room_thresh}
\end{equation}

with $\delta_T = 0.3$\,°C and $\delta_H = 0.5$\,\%.
For patient data:

\begin{equation}
	\text{Upload}_{\text{patient}} \;\Leftrightarrow\;
	\text{BPM}_n \neq \hat{\text{BPM}} \;\lor\;
	\left| \text{SpO}_{2,n} - \widehat{\text{SpO}_2} \right| > \delta_S
	\label{eq:patient_thresh}
\end{equation}

with $\delta_S = 2$\,\%.
Both conditions are evaluated every $T_{\text{interval}} = 1\,000$\,ms.

\subsubsection{End-to-End Latency Model}

The total latency from sensor event to database record is:

\begin{equation}
	L_{\text{total}} = T_{\text{acq}} + T_{\text{proc}} + T_{\text{serial}} +
	T_{\text{TLS}} + T_{\text{RTT}}
	\label{eq:latency}
\end{equation}

where:
\begin{itemize}
	\item $T_{\text{acq}}$: sensor acquisition time (DHT22: up to 2\,000\,ms;
	      MAX30102: continuous FIFO drain, $\approx$25\,ms per sample at default rate).
	\item $T_{\text{proc}}$: local computation ($<$5\,ms per SpO\textsubscript{2}
	      update; $<$1\,ms for BPM average).
	\item $T_{\text{serial}}$: JSON string build ($<$1\,ms).
	\item $T_{\text{TLS}}$: TLS handshake on first connection (1--3\,s);
	      session reuse on subsequent requests ($<$5\,ms).
	\item $T_{\text{RTT}}$: HTTPS round-trip to Firebase
	      ($\approx$100--300\,ms over 2.4\,GHz Wi-Fi).
\end{itemize}

\subsubsection{Key Assumptions}

\begin{itemize}
	\item 32-bit aligned float reads and writes are atomic on the ESP32
	      architecture, making inter-core sharing of \texttt{float} globals safe for
	      single-writer/single-reader access without a mutex.
	\item Server-side timestamps (\texttt{".sv":"timestamp"}) eliminate the need
	      for an RTC on the device.
	\item A finger is considered present when the IR raw value exceeds 50\,000 counts;
	      below this threshold BPM and SpO\textsubscript{2} are reset to zero.
\end{itemize}

\subsection{Backend Subsystem Methodology}
The backend logic is centered on a \textbf{typed procedure layer} with session-based access control. Firebase ID tokens are verified at login, then converted into a signed session cookie. Protected procedures require a valid session and enforce role-based filtering at the data layer. Alerting for presence events (unknown person, patient absence) is triggered by detection events logged from the dashboard.

\begin{figure}[H]
    \centering
	\includegraphics[width=\linewidth]{Face Recognition Pipeline.png}
    \caption{Face recognition pipeline and clinical alert decision flow.}
    \label{fig:alert_logic}
\end{figure}

\subsection{Web Dashboard Subsystem Methodology}
The dashboard utilizes a \textbf{reactive component architecture} built with React Hooks. It maintains a real-time connection to the Firebase Realtime Database using a WebSocket-based subscription model. For security and presence tracking, the methodology includes local inference using \textbf{face-api.js} to detect and identify authorized medical personnel. Visualizations are rendered using a rolling buffer of vitals, updated every 1,000\,ms to maintain a high-fidelity trend view. A camera detection loop runs every 600\,ms; alerts are issued only after a stability threshold of 5 consecutive frames.

\subsection{AI Subsystem Methodology}
The AI subsystem formulates the monitoring task as a \textbf{time-series anomaly detection and forecasting problem}.
\begin{itemize}
		\item \textbf{Anomaly Detection:} Calculated via the Z-score $z = \frac{x - \mu}{\sigma}$ over a sliding window of the most recent samples (target length 20, minimum 5). Values exceeding a magnitude of 2.5 are flagged as outliers.
		\item \textbf{Forecasting:} Implemented via \textbf{Linear Trend Estimation}, where the next 5 samples are projected using the slope $m$ of the current window.
		\item \textbf{Risk Assessment:} A rule-based classifier synthesizes the anomalies and trends into a discrete clinical status (stable, warning, critical).
\end{itemize}

\subsection{Face Recognition Methodology}

The face recognition pipeline executes fully in the browser, using a webcam stream
and pre-trained face-api.js models (TinyFaceDetector, FaceLandmark68Net,
FaceRecognitionNet, FaceExpressionNet). Each detected face produces a descriptor
vector $\mathbf{d} \in \mathbb{R}^{128}$. Matching uses Euclidean distance with a
confidence threshold of 0.6 (equivalent to a distance below 0.4):

\begin{equation}
\text{confidence} = 1 - \lVert \mathbf{d}_{\text{det}} - \mathbf{d}_{\text{enrolled}} \rVert_2
\label{eq:face_confidence}
\end{equation}

\subsection{Alert Stabilization and Rate Limiting}

The dashboard applies a multi-frame stability gate to reduce false positives.
Let $N_s = 5$ be the required number of consecutive frames. Alerts for unknown
persons and patient absence are triggered only when the condition persists for
$N_s$ frames. An additional cooldown of 30--120\,s prevents repeated alerts for
the same entity, and SMTP emails are rate limited with a 15-minute cooldown.

\subsection{Algorithmic Pseudocode}

\begin{algorithm}[H]
\caption{ESP32 sensor update loop (Core 1)}
\begin{algorithmic}[1]
\STATE Every loop iteration:
\IF{time since last DHT read $\geq 2\,s$}
	\STATE read DHT22; update $T$, $H$ if valid
\ENDIF
\STATE drain MAX30102 FIFO; if $IR < 50{,}000$ reset BPM/SpO$_2$
\STATE compute BPM from peak detection; update rolling mean
\STATE compute SpO$_2$ from 100-sample buffer; shift buffer by 25 samples
\end{algorithmic}
\end{algorithm}

\begin{algorithm}[H]
\caption{Firebase upload task (Core 0)}
\begin{algorithmic}[1]
\STATE wait for Wi-Fi; initialize Firebase with email/password
\LOOP
	\STATE app.loop() to process auth events
	\IF{Wi-Fi connected \AND app ready}
		\IF{elapsed $\geq 1\,s$}
			\STATE push room record if $|T-\hat{T}|>0.3$ or $|H-\hat{H}|>0.5$
			\STATE push patient record if BPM changed or $|SpO_2-\hat{S}|>2$
		\ENDIF
	\ELSE
		\STATE attempt reconnect every 10\,s; update Wi-Fi scan list
	\ENDIF
\ENDLOOP
\end{algorithmic}
\end{algorithm}

\begin{algorithm}[H]
\caption{Client-side face recognition loop}
\begin{algorithmic}[1]
\STATE every 600\,ms capture frame from camera
\STATE detect faces and extract descriptors
\FOR{each descriptor}
	\STATE find best match; compute confidence
	\IF{confidence $>0.6$}
		\STATE log \texttt{person recognized} event
	\ELSE
		\STATE increment unknown counter
	\ENDIF
\ENDFOR
\STATE issue alerts only if condition persists for $N_s$ frames
\end{algorithmic}
\end{algorithm}

% ============================================================
\section{Methodology Diagram}
% ============================================================

Figure~\ref{fig:methodology} presents the overall development methodology for the
complete system, from initial design through to deployment.

\begin{figure}[H]
	\centering
	\includegraphics[width=\linewidth]{methodology_diagram.png}
	\caption{Overall development methodology.}
	\label{fig:methodology}
\end{figure}

